% GENERATED FILE. Edit canonical lessons, then run npm run generate:books.
% canonical-source-hash: fnv1a64:f849c7440cd399d9
% canonical-lessons: IT-C461-quercia, IT-C461-pino, IT-C461-abete, IT-C461-potare, IT-C461-corteccia

\chapter{Gli alberi}
\label{ch:it-gli-alberi}
\hlchaptermodality{461}

\begin{chapteropening}
\textbf{By the end of this chapter:} \emph{I can name trees and care for them: the oak, the pine, the fir, pruning and the bark.}

Complete the last lesson of chapter 461: I can name trees and care for them: the oak, the pine, the fir, pruning and the bark.
\end{chapteropening}

\section[la quercia]{la quercia — the oak}

\label{lesson:IT-C461-quercia}



\begin{quote}
Before the new one: say the Italian for the bush, then the Italian for to wade across.
\end{quote}

\subsection*{You'll want to know: la quercia}
\textbf{la quercia} — \textquotedblleft{}the oak\textquotedblright{}.

\begin{quote}
Dietro casa c'è una quercia di cent'anni.
\end{quote}

\begin{quote}
\emph{Behind the house there is a hundred-year-old oak.}
\end{quote}

\emph{Forte come una quercia} — as strong as an oak.

\begin{cousinweb}
\emph{quercia} is inherited from Latin \emph{quercea} (EtymDB), from \emph{quercus} 'oak'.
\end{cousinweb}

\begin{grammarlens}[title={What you've built}]
One more word for telling, explaining and describing at B1.
\end{grammarlens}

\subsection*{Your turn}
\begin{itemize}
\raggedright
  \item \emph{Say it:} \emph{la quercia}
  \item \emph{Say it:} \emph{la quercia}, once more
  \item \emph{Say it:} a sentence of your own with \emph{la quercia}
  \item \emph{Recall:} say the Italian for the bush, then the Italian for to wade across, then say \emph{la quercia} again
\end{itemize}

\begin{tcolorbox}[breakable,colback=teal!4,colframe=teal!35!black,title={Before you move on}]
What does la quercia mean? (\textquotedblleft{}The oak\textquotedblright{}.) Say it once more.
\end{tcolorbox}
\section[il pino]{il pino — the pine (tree)}

\label{lesson:IT-C461-pino}



\begin{quote}
Before the new one: say the Italian for to wade across, then the Italian for the oak.
\end{quote}

\subsection*{You'll want to know: il pino}
\textbf{il pino} — \textquotedblleft{}the pine (tree)\textquotedblright{}.

\begin{quote}
Sulla spiaggia c'è ombra sotto i pini.
\end{quote}

\begin{quote}
\emph{On the beach there is shade under the pines.}
\end{quote}

\emph{La pineta} — a pine wood; \emph{i pinoli} are pine nuts.

\begin{cousinweb}
\emph{pino} is inherited from Latin \emph{pīnus} 'pine' (EtymDB).
\end{cousinweb}

\begin{grammarlens}[title={What you've built}]
One more word for telling, explaining and describing at B1.
\end{grammarlens}

\subsection*{Your turn}
\begin{itemize}
\raggedright
  \item \emph{Say it:} \emph{il pino}
  \item \emph{Say it:} \emph{il pino}, once more
  \item \emph{Say it:} a sentence of your own with \emph{il pino}
  \item \emph{Recall:} say the Italian for to wade across, then the Italian for the oak, then say \emph{il pino} again
\end{itemize}

\begin{tcolorbox}[breakable,colback=teal!4,colframe=teal!35!black,title={Before you move on}]
What does il pino mean? (\textquotedblleft{}The pine (tree)\textquotedblright{}.) Say it once more.
\end{tcolorbox}
\section[l'abete]{l'abete — the fir (tree)}

\label{lesson:IT-C461-abete}



\begin{quote}
Before the new one: say the Italian for the oak, then the Italian for the pine.
\end{quote}

\subsection*{You'll want to know: l'abete}
\textbf{l'abete} — \textquotedblleft{}the fir (tree)\textquotedblright{}.

\begin{quote}
In montagna ci sono boschi di abeti.
\end{quote}

\begin{quote}
\emph{In the mountains there are fir forests.}
\end{quote}

The usual Christmas tree.

\begin{cousinweb}
\emph{abete} is inherited from Latin \emph{abiēs, abiētem} 'fir' (EtymDB).
\end{cousinweb}

\begin{grammarlens}[title={What you've built}]
One more word for telling, explaining and describing at B1.
\end{grammarlens}

\subsection*{Your turn}
\begin{itemize}
\raggedright
  \item \emph{Say it:} \emph{l'abete}
  \item \emph{Say it:} \emph{l'abete}, once more
  \item \emph{Say it:} a sentence of your own with \emph{l'abete}
  \item \emph{Recall:} say the Italian for the oak, then the Italian for the pine, then say \emph{l'abete} again
\end{itemize}

\begin{tcolorbox}[breakable,colback=teal!4,colframe=teal!35!black,title={Before you move on}]
What does l'abete mean? (\textquotedblleft{}The fir (tree)\textquotedblright{}.) Say it once more.
\end{tcolorbox}
\section[potare]{potare — to prune}

\label{lesson:IT-C461-potare}



\begin{quote}
Before the new one: say the Italian for the pine, then the Italian for the fir.
\end{quote}

\subsection*{You'll want to know: potare}
\textbf{potare} — \textquotedblleft{}to prune\textquotedblright{}.

\begin{quote}
In inverno il nonno pota gli alberi.
\end{quote}

\begin{quote}
\emph{In winter Grandad prunes the trees.}
\end{quote}

\emph{La potatura} — pruning.

\begin{cousinweb}
\emph{potare} is inherited from Latin \emph{putāre} 'to prune, to cut back' (EtymDB), the same verb that also meant to reckon.
\end{cousinweb}

\begin{grammarlens}[title={What you've built}]
One more word for telling, explaining and describing at B1.
\end{grammarlens}

\subsection*{Your turn}
\begin{itemize}
\raggedright
  \item \emph{Say it:} \emph{potare}
  \item \emph{Say it:} \emph{potare}, once more
  \item \emph{Say it:} a sentence of your own with \emph{potare}
  \item \emph{Recall:} say the Italian for the pine, then the Italian for the fir, then say \emph{potare} again
\end{itemize}

\begin{tcolorbox}[breakable,colback=teal!4,colframe=teal!35!black,title={Before you move on}]
What does potare mean? (\textquotedblleft{}To prune\textquotedblright{}.) Say it once more.
\end{tcolorbox}
\section[la corteccia]{la corteccia — the bark (of a tree)}

\label{lesson:IT-C461-corteccia}



\begin{quote}
Before the new one: say the Italian for the fir, then the Italian for to prune.
\end{quote}

\subsection*{You'll want to know: la corteccia}
\textbf{la corteccia} — \textquotedblleft{}the bark (of a tree)\textquotedblright{}.

\begin{quote}
La corteccia del pino è ruvida.
\end{quote}

\begin{quote}
\emph{Pine bark is rough.}
\end{quote}

Also the cortex of the brain: \emph{la corteccia cerebrale}.

\begin{cousinweb}
\emph{corteccia} is inherited from Latin \emph{corticea} (EtymDB), from \emph{cortex, corticis} 'bark'.
\end{cousinweb}

\begin{grammarlens}[title={What you've built}]
One more word for telling, explaining and describing at B1.
\end{grammarlens}

\subsection*{Your turn}
\begin{itemize}
\raggedright
  \item \emph{Say it:} \emph{la corteccia}
  \item \emph{Say it:} \emph{la corteccia}, once more
  \item \emph{Say it:} a sentence of your own with \emph{la corteccia}
  \item \emph{Recall:} say the Italian for the fir, then the Italian for to prune, then say \emph{la corteccia} again
\end{itemize}

\begin{tcolorbox}[breakable,colback=teal!4,colframe=teal!35!black,title={Before you move on}]
What does la corteccia mean? (\textquotedblleft{}The bark (of a tree)\textquotedblright{}.) Say it once more.
\end{tcolorbox}
